\documentclass{article}

% NeurIPS 2024 style
\usepackage{neurips_2024}

\usepackage[utf8]{inputenc}
\usepackage[T1]{fontenc}
\usepackage{hyperref}
\usepackage{url}
\usepackage{booktabs}
\usepackage{amsfonts}
\usepackage{nicefrac}
\usepackage{microtype}
\usepackage{xcolor}
\usepackage{amsmath}
\usepackage{amssymb}
\usepackage{mathtools}
\usepackage{graphicx}
\usepackage{algorithm}
\usepackage{algpseudocode}
\usepackage{subcaption}
\usepackage{multirow}
\usepackage{wrapfig}

% custom math operators
\DeclareMathOperator*{\argmin}{arg\,min}
\DeclareMathOperator*{\argmax}{arg\,max}
\newcommand{\RR}{\mathbb{R}}
\newcommand{\EE}{\mathbb{E}}
\newcommand{\bx}{\mathbf{x}}
\newcommand{\bw}{\mathbf{w}}
\newcommand{\bp}{\mathbf{p}}
\newcommand{\bA}{\mathbf{A}}
\newcommand{\bs}{\mathbf{s}}

\title{ComponentMAE: Unsupervised Image Segmentation via\\
Routed Masked Autoencoders with Cross-Mask Consistency}

\author{%
  Ilias Paralikas \\
  \texttt{paralikasilias@gmail.com}
}

\begin{document}

\maketitle

\begin{abstract}
We present \textbf{ComponentMAE}, a method for unsupervised image segmentation that
combines a Masked Autoencoder (MAE) backbone with a differentiable routing network.
The key insight is that if $N$ independent decoders are forced to each explain
different spatial regions of an image, the segmentation boundaries that emerge from
this competition reflect meaningful semantic structure---without any pixel-level
annotation.

ComponentMAE employs a shared Vision Transformer (ViT) encoder whose latent
representations are routed, via soft segmentation maps produced by a UNet, to $N$
component-specific MAE decoders.  The final reconstruction is a per-pixel weighted
sum of component reconstructions, where the weights are the segmentation probabilities.
Beyond the standard masked reconstruction objective, we introduce three novel
training signals: (i) a \emph{component reconstruction loss} that assigns equal
gradient to every decoder regardless of routing, preventing component collapse;
(ii) a \emph{cross-mask consistency loss} (CM loss) that exploits complementary
two-pass masking to penalise each decoder for failing to reconstruct patches it
attended to in the complementary pass, with an area-normalisation that gives
equal gradient to semantically small components; and (iii) a \emph{routing loss}
that uses decoder reconstruction quality as a hard pseudo-label to train the
segmentation UNet.  These signals are complemented by adversarial sharpening
through a patch-level Wasserstein GAN with gradient penalty and feature-matching
loss.  Together, these objectives induce segmentation maps that capture coherent,
spatially contiguous semantic regions in a fully unsupervised manner.
\end{abstract}

% ─────────────────────────────────────────────────────────────────────────────
\section{Introduction}
\label{sec:intro}
% ─────────────────────────────────────────────────────────────────────────────

Image segmentation---partitioning an image into semantically coherent regions---is
a cornerstone of computer vision.  Despite decades of progress, supervised
segmentation methods remain expensive to apply to new domains because they require
dense pixel-level annotations, which cost orders of magnitude more to collect than
image-level labels.  This has motivated a sustained research interest in
\emph{unsupervised} segmentation: discovering semantic structure without human
labels.

Recent self-supervised learning methods have dramatically closed the gap to
supervised performance on downstream tasks.  Among these, Masked Autoencoders
(MAE)~\cite{he2022masked} have emerged as a particularly powerful pre-training
strategy: by masking a large fraction of image patches and training a ViT encoder--decoder
to reconstruct the missing content, MAE forces the encoder to develop holistic
semantic representations.  Crucially, this masking-and-reconstruction objective
creates an information bottleneck---the encoder must distil global context into a
compact latent that is then used by the decoder to ``hallucinate'' masked content.
This property makes MAE an attractive backbone for segmentation: a decoder that
successfully reconstructs a masked patch must be reasoning about the object to
which that patch belongs.

We exploit this property directly.  Instead of using MAE purely for pre-training
a feature extractor, we \emph{integrate segmentation into the MAE training loop
itself}.  ComponentMAE replaces the single MAE decoder with $N$ component-specific
decoders whose contributions to the final reconstruction are weighted by a learned
segmentation network.  If component decoder $n$ is responsible for reconstructing
a car, it must receive car-related context from the encoder; if it is also asked to
reconstruct road pixels, it faces a conflicting objective.  This competition drives
the segmentation maps toward semantically consistent boundaries.

The core technical challenges we address are:
\begin{enumerate}
  \item \textbf{Component collapse.}  With $N$ decoders, the trivially easy
    solution is for one decoder to reconstruct everything and for the routing
    network to assign weight~1 to it.  We prevent this via a dedicated
    \emph{component reconstruction loss} that gives every decoder equal gradient
    regardless of the routing.

  \item \textbf{Gradient imbalance across component sizes.}  Semantic classes
    cover very different fractions of the image---in urban street scenes, road
    and sky dominate over pedestrians.  Naive losses give large-area classes
    proportionally larger gradient, causing small classes to vanish.  We design
    the \emph{cross-mask consistency loss} with an explicit area normalisation
    that equalises the gradient contribution per component.

  \item \textbf{Spatially smooth, sharp segmentation.}  Reconstruction objectives
    alone tend to produce blurry, per-patch decisions.  We add total variation
    regularisation on the segmentation maps and adversarial sharpening of the
    reconstructed patches via a patch-level Wasserstein discriminator.

  \item \textbf{Complementary information across passes.}  Standard MAE uses a
    single masked pass, leaving exactly the visible patches unpenalised.  We
    introduce a two-pass masking strategy with complementary masks so that
    (a) every patch is reconstructed in exactly one of the two passes, and
    (b) the attention maps from pass A can be used to construct a consistency
    signal in pass B---the cross-mask consistency loss.
\end{enumerate}

Our contributions are:
\begin{itemize}
  \item The ComponentMAE architecture: a shared ViT encoder with $N$ parallel
    MAE decoders routed by a UNet segmentation network.
  \item The two-pass masking strategy enabling full-image gradient coverage.
  \item The cross-mask consistency loss with area-independent weighting, derived
    from decoder self-attention maps.
  \item A component reconstruction (anti-collapse) loss with masked-patch
    restriction.
  \item A reconstruction-guided routing loss providing pseudo-labels from
    decoder quality.
  \item A patch-level WGAN-GP with feature matching for perceptual sharpness.
\end{itemize}

% ─────────────────────────────────────────────────────────────────────────────
\section{Related Work}
\label{sec:related}
% ─────────────────────────────────────────────────────────────────────────────

\paragraph{Masked Autoencoders.}
He~et~al.~\cite{he2022masked} introduced MAE, masking 75\% of image patches and
training a ViT encoder--decoder to reconstruct raw pixel values in the masked
positions.  The asymmetric design (heavy encoder, lightweight decoder) encourages
the encoder to develop rich semantic representations.  BEiT~\cite{bao2022beit}
similarly uses a masked prediction objective but predicts discrete visual tokens
rather than raw pixels.  SimMIM~\cite{xie2022simmim} and MaskFeat~\cite{wei2022masked}
explore alternative reconstruction targets including HOG features.  ComponentMAE
differs from all of these in that it augments the MAE objective with a segmentation
routing mechanism trained jointly with the reconstruction.

\paragraph{DINO and self-distillation.}
DINO~\cite{caron2021emerging} demonstrated that Vision Transformers trained with
self-distillation produce attention maps that segment foreground objects without
supervision.  DINOv2~\cite{oquab2024dinov2} scales this to large curated datasets.
iBot~\cite{zhou2022image} combines masked image modelling with self-distillation.
These methods require a teacher network and rely on emergent attention properties
rather than explicitly optimising for segmentation.  ComponentMAE directly
supervises segmentation maps through reconstruction quality.

\paragraph{Unsupervised segmentation.}
STEGO~\cite{hamilton2022unsupervised} propagates DINO features through a
self-supervised correspondence loss to obtain dense segmentation.
MaskContrast~\cite{van2021unsupervised} uses contrastive learning on saliency-guided
superpixels.  PiCIE~\cite{cho2021picie} learns segmentation by enforcing
photometric and geometric invariances.  HP~\cite{ke2022unsupervised} hierarchically
partitions features.  These approaches treat segmentation as post-processing on
pre-trained features; ComponentMAE instead \emph{co-trains} segmentation with
reconstruction.  Closest in spirit to our work is the slot-attention
line~\cite{locatello2020object,seitzer2022bridging}: Slot Attention decomposes
scenes into $K$ learned ``slots'' through iterative attention, and DINOSAUR
\cite{seitzer2022bridging} combines slots with DINO features for object-centric
representation.  Our routing mechanism can be viewed as a non-iterative alternative
to slot attention that integrates naturally with the MAE training loop.

\paragraph{Segment Anything Model (SAM).}
SAM~\cite{kirillov2023segment} trains a prompt-conditioned segmentation model on
a massive curated dataset of 1B masks.  While powerful, SAM is not unsupervised---it
requires extensive human annotation in its training set and relies on user prompts
at inference.  ComponentMAE targets the regime where no annotations are available.

\paragraph{GAN-based segmentation.}
DatasetGAN~\cite{zhang2021datasetgan} and SemanticGAN~\cite{li2021semantic} extract
segmentation from the internal activations of a pre-trained GAN generator.
Barbershop~\cite{zhu2021barbershop} and related work exploit GAN latent spaces for
semantic manipulation.  These methods are limited to the distribution captured by
the GAN and require either the GAN training set or extensive latent optimisation.
We use an adversarial discriminator only for perceptual sharpening, not as the
primary segmentation mechanism.

\paragraph{Feature matching in GANs.}
Feature matching was introduced by Salimans~et~al.~\cite{salimans2016improved} as
a training stabiliser.  Pix2pix~\cite{isola2017image} popularised perceptual feature
matching in conditional image synthesis.  We apply feature matching specifically to
masked patch reconstruction, weighting by cross-mask consistency importance weights,
providing a content-aware perceptual loss that is more stable than the minimax
objective alone.

% ─────────────────────────────────────────────────────────────────────────────
\section{Method}
\label{sec:method}
% ─────────────────────────────────────────────────────────────────────────────

\subsection{Architecture Overview}
\label{sec:arch}

\begin{figure}[t]
  \centering
  % Placeholder — replace with actual architecture diagram
  \fbox{\parbox{0.95\textwidth}{\centering\vspace{3cm}
  \textbf{[Architecture Figure Placeholder]}\\[4pt]
  Left: shared ViT encoder processes visible patches from the two-pass masked
  input.  Center: $N$ independent MAE decoders (one per component) produce
  per-component patch reconstructions.  Right: UNet segmentation network
  produces soft routing maps; the blended image is a weighted sum of
  component images.  Bottom: loss terms and their connections to model
  components.
  \vspace{3cm}}}
  \caption{ComponentMAE architecture.  A single ViT encoder is shared across all
  $N$ component decoders.  The UNet segmentation network operates on the original
  image and produces component probability maps that blend the decoder outputs into
  the final reconstruction.  All modules are trained jointly via a composite loss.}
  \label{fig:arch}
\end{figure}

ComponentMAE consists of three trainable modules (Figure~\ref{fig:arch}):

\paragraph{Shared ViT encoder.}
The encoder $f_\theta$ follows the standard MAE design~\cite{he2022masked}: a
patch embedding layer projects each $P \times P$ image patch into a
$d_\mathrm{enc}$-dimensional token; fixed 2D sinusoidal positional embeddings
are added; then $L_\mathrm{enc}$ pre-LN transformer blocks process the
\emph{visible} tokens only.  No CLS token is used.  We use
$P\!=\!16$, $d_\mathrm{enc}\!=\!768$, $L_\mathrm{enc}\!=\!12$, $H_\mathrm{enc}\!=\!12$ heads
for a total encoder depth equal to ViT-Base.

Formally, given image $\bx \in \RR^{C \times H \times W}$ with $N_p = (H/P)^2$
patches, the encoder produces:
\begin{equation}
  \mathbf{z} = f_\theta\!\left(\{\bx_i\}_{i \notin \mathcal{M}}\right)
  \in \RR^{N_\mathrm{vis} \times d_\mathrm{enc}},
\end{equation}
where $\mathcal{M} \subset [N_p]$ is the set of masked patch indices and
$N_\mathrm{vis} = N_p - |\mathcal{M}|$.

\paragraph{Component decoders.}
$N$ independent, parameter-disjoint MAE decoders
$\{g_{\phi_n}\}_{n=1}^{N}$ each receive the same encoder latent $\mathbf{z}$.
Each decoder: (1) projects $\mathbf{z}$ to the decoder dimension
$d_\mathrm{dec} = 512$ via a linear layer; (2) inserts a learned
mask token at each masked position and un-shuffles to the original patch
order; (3) adds decoder positional embeddings; (4) runs
$L_\mathrm{dec}$ pre-LN transformer blocks; (5) applies a linear prediction
head to produce raw patch predictions $\hat{\bp}_n \in \RR^{N_p \times (C P^2)}$.
We use $L_\mathrm{dec}\!=\!4$-$12$, $d_\mathrm{dec}\!=\!512$, $H_\mathrm{dec}\!=\!16$ heads.

The image-space reconstruction for component $n$ is obtained by applying
$\tanh$ and unpatchifying:
\begin{equation}
  \hat{\bx}_n = \mathrm{unpatchify}\!\left(\tanh(\hat{\bp}_n)\right)
  \in \RR^{C \times H \times W}.
\end{equation}

\paragraph{Segmentation UNet.}
A standard UNet $s_\psi$ with skip connections takes the full image as input
and produces $N$ channel logits, which are passed through a channel-wise softmax:
\begin{equation}
  \mathbf{S} = \mathrm{softmax}(s_\psi(\bx)) \in \RR^{N \times H \times W},
  \quad \sum_{n=1}^{N} S_n(r,c) = 1 \;\forall\, (r,c).
\end{equation}
$S_n(r,c)$ is the probability that pixel $(r,c)$ belongs to component $n$.

\paragraph{Blended reconstruction.}
The final reconstruction is the spatially-weighted sum:
\begin{equation}
  \hat{\bx} = \sum_{n=1}^{N} S_n \odot \hat{\bx}_n,
  \label{eq:blend}
\end{equation}
where $\odot$ denotes element-wise multiplication with broadcasting over the channel dimension.

\subsection{Two-Pass Masking Strategy}
\label{sec:twopass}

Standard MAE uses a single random mask, leaving visible patches completely
unsupervised by the reconstruction objective.  We instead draw a single random
binary mask $\mathbf{m}^A \in \{0,1\}^{N_p}$ with mask ratio $r\!=\!0.5$, and
define the complementary mask $\mathbf{m}^B = \mathbf{1} - \mathbf{m}^A$.
With a 50\% mask ratio, $\mathbf{m}^B$ has the same cardinality as $\mathbf{m}^A$,
and their union covers \emph{all} patches:
$\mathcal{M}_A \cup \mathcal{M}_B = [N_p]$,\ \  $\mathcal{M}_A \cap \mathcal{M}_B = \emptyset$.

We run two forward passes through the encoder and all $N$ decoders:
\begin{align}
  \text{Pass A:} \quad &\hat{\bp}_n^A,\; \bA_n^A = g_{\phi_n}\!\left(f_\theta(\bx, \mathbf{m}^A),\; \mathbf{m}^A,\; \text{return\_attn=True}\right), \\
  \text{Pass B:} \quad &\hat{\bp}_n^B = g_{\phi_n}\!\left(f_\theta(\bx, \mathbf{m}^B),\; \mathbf{m}^B\right),
\end{align}
where $\bA_n^A \in \RR^{N_p \times N_p}$ is the head-averaged self-attention matrix
from the last transformer block of decoder $n$ during pass A.

This design has two benefits.  First, every patch receives a reconstruction
gradient in at least one pass.  Second, the visible patches in pass A become the
\emph{reconstruction targets} in pass B, enabling the cross-mask consistency loss
described in Section~\ref{sec:cmloss}.

\subsection{Masked Reconstruction and Component Reconstruction Losses}
\label{sec:reconloss}

\paragraph{Masked reconstruction loss ($\mathcal{L}_\mathrm{rec}$).}
The primary training signal is an MSE between the blended reconstruction
(Eq.~\ref{eq:blend}) and the target image, restricted to masked patches only.
Restricting to masked patches is essential: visible patches have their exact
pixel values directly available to the encoder, so reconstruction there is
trivially easy and provides no useful learning signal.

\begin{equation}
  \mathcal{L}_\mathrm{rec} = \frac{1}{|\mathcal{M}|} \sum_{i \in \mathcal{M}}
  \left\|\hat{\bx}_{[i]} - \bx_{[i]}\right\|^2,
\end{equation}
where $\hat{\bx}_{[i]}$ and $\bx_{[i]}$ denote the pixels belonging to patch $i$.

\paragraph{Component reconstruction loss ($\mathcal{L}_\mathrm{comp}$).}
With $N$ independent decoders and a routing network, the trivially optimal solution
is for one decoder to reconstruct everything while the routing network assigns weight
1 to that decoder and 0 to the rest.  We call this \emph{component collapse}, and
it prevents any useful segmentation from emerging.

To prevent collapse, we apply a uniform-weight MSE to every decoder independently,
\emph{before} blending, also restricted to masked patches:
\begin{equation}
  \mathcal{L}_\mathrm{comp} = \frac{1}{N} \sum_{n=1}^{N}
  \frac{1}{|\mathcal{M}|} \sum_{i \in \mathcal{M}}
  \left\|\hat{\bx}_{n,[i]} - \bx_{[i]}\right\|^2.
\end{equation}
Crucially, this loss does \emph{not} involve the segmentation weights $S_n$.
Every decoder receives identical gradient from $\mathcal{L}_\mathrm{comp}$
regardless of routing, ensuring that all decoders remain competent
reconstructors and the routing network is not rewarded for simply suppressing
all-but-one decoder.

\subsection{Cross-Mask Consistency Loss}
\label{sec:cmloss}

\paragraph{Intuition.}
In pass A, decoder $n$ reconstructs masked patch $j$ (e.g.\ a car window) by
attending to visible patch $i$ (e.g.\ a car door).  If the segmentation network
assigns high probability to component $n$ for patch $j$---meaning component $n$
is primarily responsible for the car region---then the information in patch $i$
must be relevant to component $n$'s semantic region.  In pass B, patch $i$ is
masked.  We therefore \emph{penalise} decoder $n$ for failing to reconstruct
patch $i$ in pass B, weighted by how much it relied on $i$ in pass A.

\paragraph{Importance weights.}
For component $n$ and visible-in-$A$ patch $i$, the importance weight is:
\begin{equation}
  w_n[i] = \sum_{j \in \mathcal{M}_A}
  \frac{s_n[j]}{\bar{s}_n} \cdot \tilde{A}_n^A[j \to i],
  \label{eq:cm_weight}
\end{equation}
where $s_n[j] = \frac{1}{P^2}\sum_{\text{pixels of }j} S_n$ is the mean
segmentation probability of component $n$ over patch $j$, $\bar{s}_n =
\frac{1}{N_p}\sum_i s_n[i]$ is the per-component mean (area fraction), and
$\tilde{A}_n^A[j \to i]$ is the attention weight from masked query $j$ to
visible key $i$, re-normalised so that attention to mask-token positions (which
carry no content) is zeroed out:
\begin{equation}
  \tilde{A}_n^A[j \to i] = \frac{A_n^A[j,i] \cdot \mathbf{1}[i \notin \mathcal{M}_A]}
  {\sum_{i'} A_n^A[j,i'] \cdot \mathbf{1}[i' \notin \mathcal{M}_A] + \epsilon}.
\end{equation}

\paragraph{Area-independent normalisation.}
The term $1/\bar{s}_n$ in Eq.~\ref{eq:cm_weight} is the key novelty.  Without it,
a component covering 50\% of the image (e.g.\ pavement) would receive ten times
more total gradient than one covering 5\% (e.g.\ pedestrians), causing the small
component to vanish or merge.  Dividing by the mean area fraction $\bar{s}_n$
normalises the total importance weight per component:
\begin{equation}
  \sum_{i} w_n[i] \approx \sum_{j \in \mathcal{M}_A} \frac{s_n[j]}{\bar{s}_n}
  \approx \frac{N_p \cdot \bar{s}_n}{\bar{s}_n} = N_p,
\end{equation}
making the expected total weight the same for all components regardless of their
spatial coverage.

\paragraph{Loss formulation.}
Given the importance weights, the CM loss is a weighted MSE between pass-B
decoder outputs (in raw patch space, before $\tanh$) and the target patches in
pass-B masked positions (which were visible in pass A):
\begin{equation}
  \mathcal{L}_\mathrm{cm} = \frac{1}{N} \sum_{n=1}^{N}
  \frac{\sum_{i} w_n[i] \cdot \left\|\hat{\bp}_n^B[i] - \bx_\mathrm{patch}[i]\right\|^2}
  {\sum_{i} w_n[i] + \epsilon}.
  \label{eq:cmloss}
\end{equation}
Pass-B decoder predictions $\hat{\bp}_n^B$ carry gradients through to the decoders
and encoder; all other quantities ($w_n$, attention weights, segmentation
probabilities) are \emph{detached} from the computational graph, so
$\mathcal{L}_\mathrm{cm}$ trains only the reconstruction model.

\paragraph{Patch-space comparison.}
We intentionally compute Eq.~\ref{eq:cmloss} in raw patch space (decoder output
before $\tanh$ and unpatchify) rather than in pixel space.  This avoids the
$\tanh \to$ unpatchify $\to$ bilinear interpolation chain, which (a) introduces
numerical rounding errors, and (b) saturates gradients through the $\tanh$
nonlinearity.  Patch-space comparison gives cleaner, unsaturated gradients directly
at the decoder prediction head.

\subsection{Routing and Regularisation Losses}
\label{sec:routing}

\paragraph{Routing loss ($\mathcal{L}_\mathrm{route}$).}
The segmentation UNet must be trained to assign high probability to the component
that best explains each pixel.  We construct hard pseudo-labels by selecting, for
each pixel, the component decoder with the lowest reconstruction MSE:
\begin{equation}
  \hat{y}(r,c) = \argmin_{n} \left(\hat{\bx}_n(r,c) - \bx(r,c)\right)^2,
\end{equation}
and train the UNet with cross-entropy loss:
\begin{equation}
  \mathcal{L}_\mathrm{route} = -\frac{1}{HW} \sum_{r,c}
  \log S_{\hat{y}(r,c)}(r,c).
\end{equation}
Component images are detached when computing the pseudo-labels, so gradients flow
only into the segmentation network $s_\psi$.

\paragraph{Entropy loss ($\mathcal{L}_\mathrm{ent}$).}
The softmax routing is susceptible to a second type of trivial solution: if the
network always outputs uniform probabilities ($S_n \equiv 1/N$ everywhere), the
blended reconstruction becomes the simple average of all component predictions, and
no segmentation information is conveyed.  We therefore maximise the entropy of the
routing distribution to encourage it to make confident, varied assignments:
\begin{equation}
  \mathcal{L}_\mathrm{ent} = \frac{1}{HW}\sum_{r,c}
  \sum_{n=1}^{N} S_n(r,c) \log S_n(r,c),
  \quad \text{(minimised $\Rightarrow$ higher entropy)}.
\end{equation}
This might appear contradictory to the cross-entropy routing loss, but the two
operate on complementary aspects: $\mathcal{L}_\mathrm{route}$ pushes toward
hard assignments at pixels where one decoder is clearly better; $\mathcal{L}_\mathrm{ent}$
prevents the entire map from collapsing to a single component.

\paragraph{Total variation loss ($\mathcal{L}_\mathrm{tv}$).}
Segmentation maps produced per-patch tend to be spatially noisy.  We add isotropic
total variation regularisation to encourage piecewise-smooth boundaries:
\begin{equation}
  \mathcal{L}_\mathrm{tv} = \frac{1}{N} \sum_{n=1}^{N}
  \sum_{r,c} \left|S_n(r\!+\!1,c) - S_n(r,c)\right|
           + \left|S_n(r,c\!+\!1) - S_n(r,c)\right|.
\end{equation}

\subsection{Adversarial Sharpness: Patch Discriminator and Feature Matching}
\label{sec:adv}

\paragraph{Patch discriminator.}
Pixel-wise MSE losses are known to produce blurry reconstructions because they
average over plausible high-frequency content~\cite{pathak2016context}.  We add
a \emph{PatchDiscriminator} that evaluates each $P \times P$ patch independently,
following the PatchGAN~\cite{isola2017image} design but operating in Wasserstein
distance~\cite{arjovsky2017wasserstein}.

The discriminator consists of three strided convolutions
($16\to 8 \to 4 \to 2$ spatial resolution) with LeakyReLU activations and no
batch normalisation (required by WGAN-GP~\cite{gulrajani2017improved}), followed
by a linear prediction head.  It outputs a scalar score per patch.  Intermediate
feature maps at each convolutional block are exposed for feature-matching loss.

\paragraph{Wasserstein critic loss.}
The discriminator (critic) is trained with:
\begin{equation}
  \mathcal{L}_D = \EE_{\tilde{\bx}}[D(\tilde{\bx})] - \EE_{\bx}[D(\bx)]
  + \lambda_\mathrm{gp}\, \EE_{\hat{\bx}}\!\left[(\|\nabla_{\hat{\bx}} D(\hat{\bx})\|_2 - 1)^2\right]
  + \lambda_\mathrm{drift}\, \EE_{\bx}[D(\bx)^2],
\end{equation}
where $\hat{\bx}$ are random convex combinations of real and fake patches
(gradient penalty), and the drift term prevents the critic output from growing
unbounded on real samples.  We use $\lambda_\mathrm{gp} = 10$ and
$\lambda_\mathrm{drift} = 0.001$.  The critic is updated $N_\mathrm{critic} = 3$
times per generator step.

\paragraph{Feature-matching generator loss.}
Rather than training the generator with the standard Wasserstein term
$-\EE_{\tilde{\bx}}[D(\tilde{\bx})]$ alone, which leads to training instability
in our setting (the generator is simultaneously optimised for many other losses),
we use a \emph{feature-matching} generator loss~\cite{salimans2016improved}:
\begin{equation}
  \mathcal{L}_\mathrm{fm} = \frac{1}{L} \sum_{\ell=1}^{L}
  \EE\!\left[\|D_\ell(\hat{\bx}) - D_\ell(\bx)\|^2\right],
\end{equation}
where $D_\ell(\cdot)$ are intermediate feature maps at layer $\ell$ of the
discriminator, and real features are computed under \texttt{no\_grad}.  The
generator is trained with:
\begin{equation}
  \mathcal{L}_G = \mathcal{L}_\mathrm{fm} + \alpha_\mathrm{rf} \cdot \mathcal{L}_\mathrm{wgan},
  \label{eq:gloss}
\end{equation}
where $\mathcal{L}_\mathrm{wgan} = -\EE_{\tilde{\bx}}[D(\tilde{\bx})]$ and
$\alpha_\mathrm{rf} = 0.1$ keeps the real/fake term as a minor contributor.

This design is deliberate.  Pure feature matching with a fixed real target avoids
the minimax game for the generator, preventing mode collapse and competing gradient
signals.  The small real/fake term $\alpha_\mathrm{rf}$ closes a ``texture shortcut''
loophole: without it, the generator could match discriminator features trivially
by blurring, since feature distances are less sensitive to high frequencies than
per-pixel MSE.  A small adversarial push forces the generator to also fool the
critic, incentivising plausible high-frequency texture without destabilising training.

\paragraph{Masked-patch restriction.}
Adversarial losses are applied \emph{only to masked patches}.  Visible patches
are copied directly from the input and need no sharpening; applying the discriminator
to them would penalise the identity mapping.  Masked patches from both passes A
and B are extracted as $K \times C \times P \times P$ tensors and fed to the
discriminator as a batch.

\subsection{Total Objective}
\label{sec:total}

The total loss splits across the three trainable modules:

\paragraph{Reconstruction model ($f_\theta, \{g_{\phi_n}\}$):}
\begin{equation}
  \mathcal{L}_\text{rec-model} =
    \lambda_\mathrm{rec}\, \mathcal{L}_\mathrm{rec}
    + \lambda_\mathrm{comp}\, \mathcal{L}_\mathrm{comp}
    + \lambda_\mathrm{cm}\,  \mathcal{L}_\mathrm{cm}
    + \lambda_\mathrm{adv}\, \mathcal{L}_G.
\end{equation}

\paragraph{Segmentation network ($s_\psi$):}
\begin{equation}
  \mathcal{L}_\text{seg} =
    \lambda_\mathrm{rec}\, \mathcal{L}_\mathrm{rec}
    + \lambda_\mathrm{route}\, \mathcal{L}_\mathrm{route}
    + \lambda_\mathrm{ent}\, \mathcal{L}_\mathrm{ent}
    + \lambda_\mathrm{tv}\, \mathcal{L}_\mathrm{tv}.
\end{equation}

\paragraph{Discriminator ($D$):}
\begin{equation}
  \mathcal{L}_\text{disc} = \mathcal{L}_D.
\end{equation}
Note that $\mathcal{L}_\mathrm{rec}$ appears in both the reconstruction model and
segmentation network objectives: the blended reconstruction (Eq.~\ref{eq:blend})
depends on both the decoder outputs and the segmentation weights, so both modules
receive gradient through it.  However, when computing $\mathcal{L}_\mathrm{route}$,
the decoder images are detached so that only the segmentation network is trained
by the pseudo-label signal.

Algorithm~\ref{alg:training} summarises one training iteration.

\begin{algorithm}[t]
\caption{ComponentMAE training iteration}
\label{alg:training}
\begin{algorithmic}[1]
\Require Batch $\{\bx_b\}_{b=1}^B$, models $f_\theta, \{g_{\phi_n}\}, s_\psi, D$
\State Sample random mask $\mathbf{m}^A$;\ \ set $\mathbf{m}^B \gets \mathbf{1} - \mathbf{m}^A$
\State $\mathbf{S} \gets \mathrm{softmax}(s_\psi(\bx))$ \Comment{segmentation probabilities}
\State \textbf{Pass A:} $\hat{\bp}_n^A, \bA_n^A \gets g_{\phi_n}(f_\theta(\bx, \mathbf{m}^A))\ \ \forall n$, return attention
\State Compute $\hat{\bx}_n \gets \mathrm{unpatchify}(\tanh(\hat{\bp}_n^A))\ \ \forall n$
\State Compute blended reconstruction $\hat{\bx} \gets \sum_n S_n \odot \hat{\bx}_n$
\State \textbf{Pass B:} $\hat{\bp}_n^B \gets g_{\phi_n}(f_\theta(\bx, \mathbf{m}^B))\ \ \forall n$
\State \textbf{Discriminator step} ($N_\mathrm{critic}$ times):
  \State \quad Extract masked patches: $\mathcal{R}^A, \mathcal{F}^A$, $\mathcal{R}^B, \mathcal{F}^B$
  \State \quad $\mathcal{L}_D \gets \mathrm{WassersteinCritic}(\mathcal{R}, \mathcal{F})$;\ \ update $D$
\State \textbf{Compute reconstruction model losses:}
  \State \quad $\mathcal{L}_\mathrm{rec} \gets \mathrm{MaskedMSE}(\hat{\bx}, \bx, \mathbf{m}^A)$
  \State \quad $\mathcal{L}_\mathrm{comp} \gets \frac{1}{N}\sum_n \mathrm{MaskedMSE}(\hat{\bx}_n, \bx, \mathbf{m}^A)$
  \State \quad $\{w_n\} \gets \mathrm{CMWeights}(\mathbf{S}, \bA_n^A, \mathbf{m}^A)$
  \State \quad $\mathcal{L}_\mathrm{cm} \gets \frac{1}{N}\sum_n \mathrm{WeightedMSE}(\hat{\bp}_n^B, \bx_\mathrm{patch}, w_n)$
  \State \quad $\mathcal{L}_G \gets \mathrm{FeatureMatching}(D, \mathcal{R}, \mathcal{F})$
  \State \quad Update $f_\theta, \{g_{\phi_n}\}$ with $\nabla(\lambda_\mathrm{rec}\mathcal{L}_\mathrm{rec} + \lambda_\mathrm{comp}\mathcal{L}_\mathrm{comp} + \lambda_\mathrm{cm}\mathcal{L}_\mathrm{cm} + \lambda_\mathrm{adv}\mathcal{L}_G)$
\State \textbf{Compute segmentation network losses:}
  \State \quad $\mathcal{L}_\mathrm{route} \gets \mathrm{CrossEntropy}(\mathbf{S}, \argmin_n \mathrm{MSE}(\hat{\bx}_n^{\text{det}}, \bx))$
  \State \quad $\mathcal{L}_\mathrm{ent} \gets \mathrm{Entropy}(\mathbf{S})$;\ \ $\mathcal{L}_\mathrm{tv} \gets \mathrm{TV}(\mathbf{S})$
  \State \quad Update $s_\psi$ with $\nabla(\lambda_\mathrm{rec}\mathcal{L}_\mathrm{rec} + \lambda_\mathrm{route}\mathcal{L}_\mathrm{route} + \lambda_\mathrm{ent}\mathcal{L}_\mathrm{ent} + \lambda_\mathrm{tv}\mathcal{L}_\mathrm{tv})$
\end{algorithmic}
\end{algorithm}

% ─────────────────────────────────────────────────────────────────────────────
\section{Training Details}
\label{sec:training}
% ─────────────────────────────────────────────────────────────────────────────

\paragraph{Dataset.}
We train on the Cityscapes dataset~\cite{cordts2016cityscapes}, which contains
19,998 training images of urban driving scenes at $256 \times 256$ resolution
after cropping and resizing.  No semantic labels are used during training;
Cityscapes annotations serve only for qualitative evaluation.

\paragraph{Architecture hyperparameters.}
The ViT encoder uses $d_\mathrm{enc}\!=\!768$, $L_\mathrm{enc}\!=\!12$,
$H_\mathrm{enc}\!=\!12$ (ViT-Base scale), $P\!=\!16$, yielding $N_p\!=\!256$ patches
per image.  Each MAE decoder uses $d_\mathrm{dec}\!=\!512$, $L_\mathrm{dec}\!=\!4$--$12$
(configuration-dependent), $H_\mathrm{dec}\!=\!16$.  The total parameter count
for a single-component model (encoder + one decoder) is approximately 124M.

The segmentation UNet uses a channel progression
$[256, 128, 64, 32, 32, 32, 16]$ with residual blocks and 4 skipped skip
connections (bottom-up path only), yielding approximately 4.2M parameters.
The PatchDiscriminator has approximately 660K parameters.

\paragraph{Optimiser and learning rate.}
All three modules use Adam optimiser.  The reconstruction model (encoder + decoders)
and segmentation UNet use learning rate $10^{-4}$.  The discriminator uses
$\text{lr}\!=\!10^{-4}$ with Adam betas $(0.0, 0.9)$ as recommended for
WGAN-GP~\cite{gulrajani2017improved}.

\paragraph{Loss weights.}
We use $\lambda_\mathrm{rec}\!=\!10$, $\lambda_\mathrm{comp}\!=\!0.1$,
$\lambda_\mathrm{cm}\!=\!0.1$ (ramped linearly from 0 over the first 10 epochs),
$\lambda_\mathrm{route}\!=\!0.01$, $\lambda_\mathrm{ent}\!=\!0.01$,
$\lambda_\mathrm{tv}\!=\!0.01$, $\lambda_\mathrm{adv}\!=\!0.01$,
$\alpha_\mathrm{rf}\!=\!0.1$.  The CM weight ramp prevents the CM loss from
interfering with early-stage reconstruction learning before the decoder
attention maps are meaningful.

\paragraph{Mask ratio.}
We use a mask ratio of $r\!=\!0.5$, lower than the typical 0.75 used in MAE
pre-training.  The lower ratio is motivated by the two-pass strategy: with
$r\!=\!0.5$, each pass sees exactly half the patches, and pass B's masked patches
are exactly pass A's visible patches.  At $r\!=\!0.75$, the passes would overlap,
reducing the CM loss signal.

\paragraph{Dropout.}
We apply dropout $p\!=\!0.2$ inside transformer attention and MLP blocks to
regularise the large encoder.

\paragraph{Training duration.}
Models are trained for 100 epochs with a batch size of 8, corresponding to
250K gradient steps.  Checkpoints of all three models are saved after each epoch.

% ─────────────────────────────────────────────────────────────────────────────
\section{Discussion}
\label{sec:discussion}
% ─────────────────────────────────────────────────────────────────────────────

\paragraph{Why MAE and not contrastive learning?}
Contrastive methods (SimCLR, MoCo, DINO) produce image-level or patch-level
representations that are excellent for classification but lack the spatial
structure needed for dense prediction: their loss functions explicitly discard
positional information via global pooling or CLS aggregation.  MAE, by contrast,
must reconstruct \emph{specific spatial locations}, forcing the representation
to retain positional and structural information.  This makes it a natural fit
for a segmentation objective that operates in the spatial domain.

\paragraph{Information flow and gradient paths.}
The gradient structure of ComponentMAE is carefully designed to avoid shortcutting.
The blended reconstruction loss $\mathcal{L}_\mathrm{rec}$ trains both the decoders
(to reconstruct well) and the segmentation network (to assign high weight to the
best-reconstructing component).  However, if we also trained the decoders via
$\mathcal{L}_\mathrm{rec}$'s gradient through the segmentation weights, the
decoders could specialise to produce outputs that are \emph{easy to blend} rather
than outputs that correctly segment.  In practice we find it beneficial to detach
the independent reconstructions when flowing gradient into the segmentation network,
so that the routing network is trained solely on the quality signal, not on decoder
specialisation artefacts.

\paragraph{The two collapses and their cures.}
ComponentMAE must resist two failure modes: (1)~\emph{component collapse}, where
one decoder handles everything --- cured by $\mathcal{L}_\mathrm{comp}$; and
(2)~\emph{entropy collapse}, where all routing weights are uniform --- cured by
$\mathcal{L}_\mathrm{ent}$ and $\mathcal{L}_\mathrm{route}$.  The tension between
these objectives acts as a force toward interesting intermediate states where
components are distinct but their spatial coverage is non-trivial.

\paragraph{Scalability.}
The parameter cost scales as $O(N)$ in the number of components, since each
component requires an independent decoder.  The attention extraction needed
for the CM loss adds one additional forward pass through the last decoder
block per component.  Both scale linearly and remain tractable for
$N \leq 16$ components on contemporary hardware.

% ─────────────────────────────────────────────────────────────────────────────
\section{Conclusion}
\label{sec:conclusion}
% ─────────────────────────────────────────────────────────────────────────────

We have presented ComponentMAE, a method for unsupervised image segmentation that
integrates a differentiable routing mechanism directly into the Masked Autoencoder
training loop.  The system comprises a shared ViT encoder, $N$ independent MAE
decoders, and a UNet segmentation network trained jointly via a suite of
complementary objectives: masked reconstruction, component anti-collapse,
cross-mask consistency with area-independent weighting, reconstruction-guided
routing, entropy regularisation, total variation, and adversarial patch sharpening.

The key insight is that if different parts of an image are better explained by
different decoders, the spatial boundaries of this explanatory advantage should
correspond to semantic boundaries.  The cross-mask consistency loss makes this
insight operational: by exploiting the attention maps that decoders produce
internally, it rewards each decoder for attending to visually consistent context
across complementary masked views---without any labels.

Future work includes: scaling to larger numbers of components and higher resolution;
integrating hierarchical segmentation (coarse-to-fine routing); replacing the
hand-tuned loss weights with a learned multi-task weighting scheme; and applying
ComponentMAE to domains beyond natural images where annotation is expensive (e.g.\
medical imaging, satellite remote sensing).

% ─────────────────────────────────────────────────────────────────────────────
\bibliographystyle{plain}
\bibliography{references}

% Inline bibliography for self-contained draft (remove when using .bib file)
\begin{thebibliography}{99}

\bibitem{he2022masked}
Kaiming He, Xinlei Chen, Saining Xie, Yanghao Li, Piotr Doll{\'a}r, and
Ross Girshick.
\newblock Masked autoencoders are scalable vision learners.
\newblock In \emph{CVPR}, 2022.

\bibitem{bao2022beit}
Hangbo Bao, Li Dong, Songhao Piao, and Furu Wei.
\newblock BEiT: BERT pre-training of image transformers.
\newblock In \emph{ICLR}, 2022.

\bibitem{xie2022simmim}
Zhenda Xie, Zheng Zhang, Yue Cao, Yutong Lin, Jianmin Bao, Zhuliang Yao,
Qi Dai, and Han Hu.
\newblock SimMIM: A simple framework for masked image modeling.
\newblock In \emph{CVPR}, 2022.

\bibitem{wei2022masked}
Chen Wei, Haoqi Fan, Saining Xie, Chao-Yuan Wu, Alan Yuille, and Christoph
Feichtenhofer.
\newblock Masked feature prediction for self-supervised visual pre-training.
\newblock In \emph{CVPR}, 2022.

\bibitem{caron2021emerging}
Mathilde Caron, Hugo Touvron, Ishan Misra, Herv{\'e} J{\'e}gou, Julien Mairal,
Piotr Bojanowski, and Armand Joulin.
\newblock Emerging properties in self-supervised vision transformers.
\newblock In \emph{ICCV}, 2021.

\bibitem{oquab2024dinov2}
Maxime Oquab, Timoth{\'e}e Darcet, Th{\'e}o Moutakanni, Huy V. Vo,
Marc Szafraniec, Vasil Khalidov, Pierre Fernandez, Daniel Haziza, Francisco
Massa, Alaaeldin El-Nouby, et~al.
\newblock DINOv2: Learning robust visual features without supervision.
\newblock \emph{TMLR}, 2024.

\bibitem{zhou2022image}
Jinghao Zhou, Chen Wei, Huiyu Wang, Wei Shen, Cihang Xie, Alan Yuille, and
Tao Kong.
\newblock iBOT: Image BERT pre-training with online tokenizer.
\newblock In \emph{ICLR}, 2022.

\bibitem{hamilton2022unsupervised}
Mark Hamilton, Zhoutong Zhang, Bharath Hariharan, Noah Snavely, and
William T. Freeman.
\newblock Unsupervised semantic segmentation by distilling feature correspondences.
\newblock In \emph{ICLR}, 2022.

\bibitem{van2021unsupervised}
Wouter Van Gansbeke, Simon Vandeghen, Maxim Berman, Marc Proesmans, and
Luc Van~Gool.
\newblock Unsupervised semantic segmentation by contrasting object mask proposals.
\newblock In \emph{ICCV}, 2021.

\bibitem{cho2021picie}
Jang Hyun Cho, Utkarsh Mall, Kavita Bala, and Bharath Hariharan.
\newblock PiCIE: Unsupervised semantic segmentation using invariance and
equivariance in clustering.
\newblock In \emph{CVPR}, 2021.

\bibitem{ke2022unsupervised}
Lei Ke, Mingqiao Ye, Martin Danelljan, Yifan Liu, Yu-Wing Tai, Chi-Keung Tang,
and Fisher Yu.
\newblock Unsupervised hierarchical semantic segmentation with multiview cosegmentation and clustering transformers.
\newblock In \emph{CVPR}, 2022.

\bibitem{locatello2020object}
Francesco Locatello, Dirk Weissenborn, Thomas Unterthiner, Aravindh Mahendran,
Georg Heigold, Jakob Uszkoreit, Alexey Dosovitskiy, and Thomas Kipf.
\newblock Object-centric learning with slot attention.
\newblock In \emph{NeurIPS}, 2020.

\bibitem{seitzer2022bridging}
Maximilian Seitzer, Max Welling, and Georg Martius.
\newblock Bridging the gap to real-world object-centric learning.
\newblock In \emph{ICLR}, 2023.

\bibitem{kirillov2023segment}
Alexander Kirillov, Eric Mintun, Nikhila Ravi, Hanzi Mao, Chloe Rolland,
Laura Gustafson, Tete Xiao, Spencer Whitehead, Alexander C. Berg,
Wan-Yen Lo, Piotr Doll{\'a}r, and Ross Girshick.
\newblock Segment anything.
\newblock In \emph{ICCV}, 2023.

\bibitem{zhang2021datasetgan}
Yuxuan Zhang, Huan Ling, Jun Gao, Kangxue Yin, Jean-Francois Lafleche,
Adela Barriuso, Antonio Torralba, and Sanja Fidler.
\newblock DatasetGAN: Efficient labeled data factory with minimal human effort.
\newblock In \emph{CVPR}, 2021.

\bibitem{li2021semantic}
Guangyao Li, Hengshuang Zhao, Yucheng Zhu, Wancheng Zhu, Lihi Zelnik-Manor,
and Luc Van~Gool.
\newblock SemanticGAN: Learning semantic segmentation from unlabeled scenes.
\newblock arXiv:2104.00841, 2021.

\bibitem{zhu2021barbershop}
Peihao Zhu, Rameen Abdal, John Femiani, and Peter Wonka.
\newblock Barbershop: GAN-based image compositing using segmentation masks.
\newblock \emph{ACM TOG}, 40(6), 2021.

\bibitem{salimans2016improved}
Tim Salimans, Ian Goodfellow, Wojciech Zaremba, Vicki Cheung, Alec Radford,
and Xi Chen.
\newblock Improved techniques for training GANs.
\newblock In \emph{NeurIPS}, 2016.

\bibitem{isola2017image}
Phillip Isola, Jun-Yan Zhu, Tinghui Zhou, and Alexei A. Efros.
\newblock Image-to-image translation with conditional adversarial networks.
\newblock In \emph{CVPR}, 2017.

\bibitem{arjovsky2017wasserstein}
Martin Arjovsky, Soumith Chintala, and L{\'e}on Bottou.
\newblock Wasserstein generative adversarial networks.
\newblock In \emph{ICML}, 2017.

\bibitem{gulrajani2017improved}
Ishaan Gulrajani, Faruk Ahmed, Martin Arjovsky, Vincent Dumoulin, and
Aaron Courville.
\newblock Improved training of Wasserstein GANs.
\newblock In \emph{NeurIPS}, 2017.

\bibitem{pathak2016context}
Deepak Pathak, Philipp Krahenbuhl, Jeff Donahue, Trevor Darrell, and Alexei A.
Efros.
\newblock Context encoders: Feature learning by inpainting.
\newblock In \emph{CVPR}, 2016.

\bibitem{cordts2016cityscapes}
Marius Cordts, Mohamed Omran, Sebastian Ramos, Timo Rehfeld, Markus Enzweiler,
Rodrigo Benenson, Uwe Franke, Stefan Roth, and Bernt Schiele.
\newblock The Cityscapes dataset for semantic urban scene understanding.
\newblock In \emph{CVPR}, 2016.

\end{thebibliography}

\end{document}
